\setcounter{exercisegroup}{0}
\refstepcounter{exercisegroup}
\section*{Exercises on \S\arabic{exercisegroup}}
\phantomsection
\addcontentsline{toc}{section}{Exercises on \protect\S\arabic{exercisegroup}}

\begin{enumerate}
\def\labelenumi{\arabic{enumi})}
\item
  Let \(\mathbf{f}\) be a continuous and differentiable map of \(\mathbf{R}\) into a complete normed algebra E over \(\mathbf{R}\) having a unit element \(\mathbf{e}\) ; suppose that \(\mathbf{f}(0) = \mathbf{e}\) and that, identically, \(\mathbf{f}'(x) = \mathbf{f}(x)\mathbf{c}\) where \(\mathbf{c}\) is an invertible element of E. Show that \(\mathbf{f}\) is a homomorphism of the additive group \(\mathbf{R}\) into the multiplicative group of invertible elements of E. (Consider the largest open interval I containing 0 such that \(\mathbf{f}(z)\) is invertible for all \(z \in I\) , and prove that for \(x, y\) and \(x + y\) in I one has \(\mathbf{f}(x + y)(\mathbf{f}(x)\mathbf{f}(y))^{-1} = \mathbf{e}\) keeping \(x\) fixed and letting \(y\) vary; then deduce that \(I = R\) .)
\item
  \begin{enumerate}
  \def\labelenumii{\alph{enumii})}
  \tightlist
  \item
    In order that the function \((1 + 1/x)^{x+p}\) should be decreasing (resp. increasing) for x \textgreater{} 0 it is necessary and sufficient that \(p \geqslant \frac{1}{2}\) (resp. \(p \leqslant 0\) ); for \(0 < p < \frac{1}{2}\) the function is decreasing on an interval {]}0, \(x_{0}[\) and increasing on {]} \(x_{0}, +\infty[\) . In every case the function tends to e as x tends to \(+\infty\) .
  \end{enumerate}
\end{enumerate}

\begin{enumerate}
\def\labelenumi{\alph{enumi})}
\setcounter{enumi}{1}
\tightlist
\item
  Study in the same way the functions \((1 - 1 / x)^{x - p}\) , \((1 + 1 / x)^{x}(1 + p / x)\) and \((1 + p / x)^{x + 1}\) for \(x > 0\) .
\end{enumerate}

\begin{enumerate}
\def\labelenumi{\arabic{enumi})}
\setcounter{enumi}{2}
\tightlist
\item
  \begin{enumerate}
  \def\labelenumii{\alph{enumii})}
  \tightlist
  \item
    Show that \((1+x)^{\alpha} \geqslant 1 + \alpha x\) for \(\alpha \geqslant 1\) and \(x \geqslant 0\) , and that \((1+x)^{\alpha} \leqslant 1 + \alpha x\) for \(0 \leqslant \alpha \leqslant 1\) and \(x \geqslant 0\) .
  \end{enumerate}
\end{enumerate}

\begin{enumerate}
\def\labelenumi{\alph{enumi})}
\setcounter{enumi}{1}
\item
  Show that \((1 - x)^{\alpha} \leqslant \frac{1}{1 + \alpha x}\) for \(0 \leqslant x \leqslant 1\) and \(\alpha \geqslant 0\) .
\item
  For \(0 \leqslant \alpha \leqslant 1\) show that for \(x \geqslant 0\) and \(y \geqslant 0\) one has \(x^{\alpha} y^{1 - \alpha} \leqslant \alpha x + by\) for every pair of numbers \(>0\) for which \(a^{\alpha} b^{1 - \alpha} = \alpha (1 - \alpha)^{1 - \alpha}\) .
\item
  Deduce from c) that if f, g are two regulated \(\geqslant 0\) functions on an interval I such that the integrals \(\int_{I} f(t) dt\) and \(\int_{I} g(t) dt\) converge, then the integral \(\int_{I} (f(t))^{\alpha}(g(t))^{1-\alpha} dt\) converges, and
\end{enumerate}

\[
\int_ {I} (f (t)) ^ {\alpha} (g (t)) ^ {1 - \alpha} d t \leqslant a \int_ {I} f (t) d t + b \int_ {I} g (t) d t.
\]

From this deduce the Hölder inequality

\[
\int_ {I} (f (t)) ^ {\alpha} (g (t)) ^ {1 - \alpha} d t \leqslant \left(\int_ {I} f (t) d t\right) ^ {\alpha} \left(\int_ {I} g (t) d t\right) ^ {1 - \alpha}
\]

(which is called the ``Cauchy-Buniakowski-Schwarz inequality'' when \(\alpha = \frac{1}{2}\) ).

\begin{enumerate}
\def\labelenumi{\arabic{enumi})}
\setcounter{enumi}{3}
\tightlist
\item
  From the Cauchy-Schwarz inequality deduce that for every function \(u\) which is a primitive of a regulated function on \([a, b]\) , and every continuous function \(h > 0\) on \([a, b]\) , one has
\end{enumerate}

\[
\left| u (b) ^ {2} - u (a) ^ {2} \right| \leqslant 2 \left(\int_ {a} ^ {b} u ^ {\prime 2} (t) h (t) d t\right) ^ {1 / 2} \left(\int_ {a} ^ {b} \frac {u ^ {2} (t)}{h (t)} d t\right) ^ {1 / 2}
\]

if the two integrals on the right-hand side converge.

By taking \(a = 0\) , \(b = \pi / 2\) , \(h(t) = 1\) and

\[
u (t) = c _ {1} \cos t + c _ {2} \cos 3 t + \dots + c _ {n} \cos (2 n - 1) t,
\]

where \(c_{j} \geqslant 0\) for every j, deduce the Carlson inequality

\[
\sum_ {j} c _ {j} \leqslant \sqrt {\pi} \left(\sum_ {j} c _ {j} ^ {2}\right) ^ {1 / 4} \left(\sum_ {j} \left(j - \frac {1}{2}\right) ^ {2} c _ {j} ^ {2}\right) ^ {1 / 4}.
\]

\begin{enumerate}
\def\labelenumi{\arabic{enumi})}
\setcounter{enumi}{4}
\tightlist
\item
  For \(x > 0\) and any real \(y\) show that
\end{enumerate}

\[
x y \leqslant \log x + e ^ {y - 1}
\]

(cf.~II, p.~82, exerc. 15); in what circumstances are the two sides equal?

\begin{enumerate}
\def\labelenumi{\arabic{enumi})}
\setcounter{enumi}{5}
\tightlist
\item
  Denote by \(A_{n}\) and \(G_{n}\) the usual arithmetic and geometric means of the first n numbers of a sequence \((a_{k})\) ( \(k \geqslant 1\) ) of positive numbers. Show that
\end{enumerate}

\[
n \left(\mathrm{A} _ {n} - \mathrm{G} _ {n}\right) \leqslant (n + 1) \left(\mathrm{A} _ {n + 1} - \mathrm{G} _ {n + 1}\right)
\]

with equality only if \(a_{n+1} = G_{n}\) (put \(a_{n+1} = x^{n+1}\) , \(G_{n} = y^{n+1}\) ).

\begin{enumerate}
\def\labelenumi{\arabic{enumi})}
\setcounter{enumi}{6}
\tightlist
\item
  With the notation of exerc. 6 show that if one puts \(x_{i} = a_{i} / \mathrm{A}_{n}\) the relation \(\mathrm{G}_n \geqslant (1 - \alpha)\mathrm{A}_n\) for \(0 \leqslant \alpha < 1\) implies, for \(1 \leqslant i \leqslant n\) ,
\end{enumerate}

\[
x _ {i} \left(1 - \frac {x _ {i} - 1}{n - 1}\right) ^ {n - 1} \geqslant (1 - \alpha) ^ {n}.
\]

Deduce that for each index \(i\) one has \(1 + x' \leqslant x_i \leqslant 1 + x''\) , where \(x'\) is the root \(\leqslant 0\) , and \(x''\) the root \(\geqslant 0\) , of the equation \((1 + x)e^{-x} = (1 - \alpha)^n\) (cf.~III, p.~115, exerc. 2).

\begin{enumerate}
\def\labelenumi{\arabic{enumi})}
\setcounter{enumi}{7}
\tightlist
\item
  Consider a sequence of sets \(D_{k}\) ( \(k \geqslant 1\) ) and for each k consider two functions \((x_{1}, \ldots, x_{k}) \mapsto f_{k}(x_{1}, \ldots, x_{k}), \quad (x_{1}, \ldots, x_{k}) \mapsto g_{k}(x_{1}, \ldots, x_{k})\) with real values, where \((x_{1}, \ldots, x_{k}) \in D_{1} \times D_{2} \times \cdots \times D_{k}\) . Suppose that for every k there is a real function \(F_{k}\) defined on R such that for every \(\mu \in R\) and arbitrary \(a_{1} \in D_{1}, \ldots, a_{k-1} \in D_{k-1}\) one has
\end{enumerate}

\[
\sup _ {x _ {k} \in \mathrm{D} _ {k}} \left(\mu f _ {k} (a _ {1}, \dots , a _ {k - 1}, x _ {k}) - g _ {k} (a _ {1}, \dots , a _ {k - 1}, x _ {k})\right) = \mathrm{F} _ {k} (\mu) f _ {k - 1} (a _ {1}, \dots , a _ {k - 1}).
\]

(By convention one takes \(f_0 = 1\) .) Show if the sequences of real numbers \((\mu_k)\) and \((\delta_k)\) satisfy the conditions \(\mathrm{F}_1(\mu_1) = 0\) and \(\mathrm{F}_k(\mu_k) = \delta_k\) then one has the inequalities

\[
\left(\mu_ {1} - \delta_ {2}\right) f _ {1} + \left(\mu_ {2} - \delta_ {3}\right) f _ {2} + \dots + \left(\mu_ {n - 1} - \delta_ {n}\right) f _ {n - 1} + \mu_ {n} f _ {n} \leqslant g _ {1} + g _ {2} + \dots + g _ {n}
\]

for all values of n and of \(x_{k} \in D_{k}\) .

\begin{enumerate}
\def\labelenumi{\arabic{enumi})}
\setcounter{enumi}{8}
\tightlist
\item
  With the notation of exerc. 6 show that
\end{enumerate}

\[
\mu_ {1} \mathrm{G} _ {1} + \mu_ {2} \mathrm{G} _ {2} + \dots + \mu_ {n} \mathrm{G} _ {n} \leqslant \lambda_ {1} a _ {1} + \lambda_ {2} a _ {2} + \dots + \lambda_ {n} a _ {n}
\]

provided that \(\lambda_{k}\geqslant 0\) for each \(k\) and

\[
\mu_ {k} = k \left(\left(\lambda_ {k} \beta_ {k}\right) ^ {1 / k} - \beta_ {k + 1} ^ {1 / k}\right) \quad \text { for } 1 \leqslant k \leqslant n
\]

with \(\beta_{k}\geqslant 0\) for all \(k\) , \(\beta_{1}\leqslant 1\) and \(\beta_{n + 1} = 0\) (method of exerc. 8).

In particular, for suitable choices of \(\lambda_{k}\) and \(\beta_{k}\) one has

\[
\mathrm{G} _ {1} + \mathrm{G} _ {2} + \dots + \mathrm{G} _ {n} + n \mathrm{G} _ {n} \leqslant 2 a _ {1} + \left(\frac {3}{2}\right) ^ {2} a _ {2} + \dots + \left(\frac {n + 1}{n}\right) ^ {n} a _ {n}
\]

and in particular (Carleman's inequality)

\[
\begin{array}{c} \mathrm{G} _ {1} + \mathrm{G} _ {2} + \dots + \mathrm{G} _ {n} <   e (a _ {1} + a _ {2} + \dots + a _ {n}); \\ e \mathrm{A} _ {n} \geqslant \Gamma_ {n} e ^ {\mathrm{G} _ {n} / \Gamma_ {n}} \end{array}
\]

where \(\Gamma_{n}=(\mathrm{G}_{1}+\mathrm{G}_{2}+\cdots+\mathrm{G}_{n})/n\) (strengthening of the Carleman inequality);

\[
\frac {\mathrm{A} _ {n}}{\mathrm{G} _ {n}} \geqslant \frac {1}{2} \left(\frac {\Gamma_ {n}}{\mathrm{G} _ {n}} + \frac {\mathrm{G} _ {n}}{\Gamma_ {n}}\right)
\]

if the sequence \((a_{k})\) is increasing;

\[
n \mathrm{G} _ {n} - m \mathrm{G} _ {m} \leqslant a _ {m + 1} + a _ {m + 2} + \dots + a _ {n}.
\]

\begin{enumerate}
\def\labelenumi{\arabic{enumi})}
\setcounter{enumi}{9}
\tightlist
\item
  The sequence \((a_{k})_{k\geqslant1}\) being formed of numbers \textgreater0, we put \(s_{n}=nA_{n}=a_{1}+\cdots+a_{n}\) . For p\textless1 and \(\lambda_{k}>0\) one has
\end{enumerate}

\[
\mu_ {1} s _ {1} ^ {1 / p} + \mu_ {2} s _ {2} ^ {1 / p} + \dots + \mu_ {n} s _ {n} ^ {1 / p} \leqslant \lambda_ {1} a _ {1} ^ {1 / p} + \lambda_ {2} a _ {2} ^ {1 / p} + \dots + \lambda_ {n} a _ {n} ^ {1 / p}
\]

provided that \(\mu_1 \geqslant \lambda_1\) and, for \(k \geqslant 2\) , \(\mu_k = \left( \lambda_k^q + \delta_k^q \right)^{1/\delta} - \delta_{k+1}\) , with \(q = p/(1-p)\) , \(\delta_k \geqslant 0\) for every \(k\) , and \(\delta_{n+1} = 0\) (method of exerc. 8). In particular, one has the inequalities

\[
\sum_ {k = 1} ^ {n} \left((k - p) ^ {1 / q} - k ^ {1 / q}\right) s _ {k} ^ {1 / q} + n ^ {1 / q} s _ {n} ^ {1 / q} \leqslant (1 - p) ^ {1 / q} \sum_ {k = 1} ^ {n} a _ {k} ^ {1 / q}
\]

which imply (with \(\mathrm{A}_n = s_n / n\) )

\[
\sum_ {k = 1} ^ {n} \mathrm{A} _ {k} ^ {1 / p} + \frac {n}{1 - p} \mathrm{A} _ {n} ^ {1 / p} \leqslant (1 - p) ^ {- 1 / p} \sum_ {k = 1} ^ {n} a _ {k} ^ {1 / p}
\]

(strengthened Hardy inequality);

\[
\frac {1}{q} \sum_ {k = 1} ^ {n} \left(b _ {k} - 1\right) A _ {k} ^ {1 / p} + n A _ {n} ^ {1 / p} \leqslant \sum_ {k = 1} ^ {n} a _ {k} ^ {1 / p} b _ {k} ^ {1 / p}
\]

provided that \(b_{k} \geqslant 1\) for all \(k\) ;

\[
\alpha \left(a _ {1} ^ {2} + a _ {2} ^ {2} + \dots + a _ {n} ^ {2}\right) \leqslant a _ {1} ^ {2} + \left(a _ {2} - a _ {1}\right) ^ {2} + \dots + \left(a _ {n} - a _ {n - 1}\right) ^ {2} + \beta a _ {n} ^ {2}
\]

if \(\alpha = 2(1 - \cos \theta)\) and \(\beta = 1 - \frac{\sin(n + 1)\theta}{\sin n\theta}\) for some \(\theta\) such that \(0\leqslant \theta < \frac{\pi}{n}\) .

Case when \(\theta = \frac{\pi}{n + 1}\) ; case when \(\theta = \frac{\pi}{2n + 1}\) .

\begin{enumerate}
\def\labelenumi{\arabic{enumi})}
\setcounter{enumi}{10}
\tightlist
\item
  Let \(a_{ij}\) ( \(1 \leqslant i \leqslant n\) , \(1 \leqslant j \leqslant n\) ) be numbers \(\geqslant 0\) such that for each index \(i\) one has \(\sum_{j=1}^{n} a_{ij} = 1\) and, for each index \(j\) one has \(\sum_{i=1}^{n} a_{ij} = 1\) . Let \(x_i\) ( \(1 \leqslant i \leqslant n\) ) be \(n\) numbers
\end{enumerate}

\[
> 0;
\]

\[
y _ {i} = a _ {i 1} x _ {1} + a _ {i 2} x _ {2} + \dots + a _ {i n} x _ {n} \quad (1 \leqslant i \leqslant n).
\]

Show that \(y_{1}y_{2}\ldots y_{n}\geqslant x_{1}x_{2}\ldots x_{n}\) (bound \(\log y_{i}\) below for each index i).

\begin{enumerate}
\def\labelenumi{\arabic{enumi})}
\setcounter{enumi}{11}
\tightlist
\item
  \begin{enumerate}
  \def\labelenumii{\alph{enumii})}
  \tightlist
  \item
    Let \(\sum_{i,j} c_{ij} x_i \bar{x}_j\) where \(c_{ji} = \bar{c}_{ij}\) be a nondegenerate positive hermitian form, and
  \end{enumerate}
\end{enumerate}

\(\Delta\) its determinant; show that \(\Delta \leqslant c_{11}c_{22}\ldots c_{nn}\) (express \(\Delta\) and the \(c_{ii}\) in terms of the eigenvalues of the hermitian form, and use prop. 2).

\begin{enumerate}
\def\labelenumi{\alph{enumi})}
\setcounter{enumi}{1}
\tightlist
\item
  Deduce that if \((a_{ij})\) is a square matrix of order n with arbitrary complex elements and \(\Delta\) is its determinant, one has (``Hadamard's inequality'')
\end{enumerate}

\[
\left| \Delta \right| ^ {2} \leqslant \left(\sum_ {j = 1} ^ {n} \left| a _ {1 j} \right| ^ {2}\right) \left(\sum_ {j = 1} ^ {n} \left| a _ {2 j} \right| ^ {2}\right) \dots \left(\sum_ {j = 1} ^ {n} \left| a _ {n j} \right| ^ {2}\right)
\]

with equality only if one of the terms on the right-hand side is zero or if, for all distinct indices \(h, k\)

\[
a _ {h 1} \bar {a} _ {k 1} + a _ {h 2} \bar {a} _ {k 2} + \dots + a _ {h n} \bar {a} _ {k n} = 0
\]

(multiply the matrix \((a_{ij})\) by the conjugate of its transpose).

\begin{enumerate}
\def\labelenumi{\arabic{enumi})}
\setcounter{enumi}{12}
\tightlist
\item
  If \(x, y, a, b\) are \(> 0\) , show that
\end{enumerate}

\[
x \log \frac {x}{a} + y \log \frac {y}{b} \geqslant (x + y) \log \frac {x + y}{a + b}
\]

with equality only if x/a = y/b.

\begin{enumerate}
\def\labelenumi{\arabic{enumi})}
\setcounter{enumi}{13}
\tightlist
\item
  Let \(a\) be a real number such that \(0 < a < \pi / 2\) . Show that the function
\end{enumerate}

\[
\frac {\frac {\tan x}{x} - \frac {\tan a}{a}}{x \tan x - a \tan a}
\]

is strictly increasing on the interval \(a, \pi/2\) (cf.~I, p.~39, exerc. 11).

\begin{enumerate}
\def\labelenumi{\arabic{enumi})}
\setcounter{enumi}{14}
\item
  Let \(u\) and \(v\) be two polynomials in \(x\) , with real coefficients, such that \(\sqrt{1 - u^2} = v\sqrt{1 - x^2}\) identically; show that, if \(n\) is the degree of \(u\) , then \(u' = nv\) ; deduce that \(u(x) = \cos(n\operatorname{Arc}\cos x)\) .
\item
  Show (by induction on \(n\) ) that
\end{enumerate}

\[
\mathrm{D} ^ {n} \left(\operatorname{Arc} \tan x\right) = (- 1) ^ {n - 1} \frac {(n - 1) !}{(1 + x ^ {2}) ^ {n / 2}} \sin \left(n \operatorname{Arc} \tan \frac {1}{x}\right).
\]

¶ 17) Let \(f\) be a real function defined on an open interval \(\mathrm{I}\subset \mathbf{R}\) , and such that for every system of three points \(x_{1}, x_{2}, x_{3}\) of I satisfying \(x_{1} < x_{2} < x_{3} < x_{1} + \pi\) one has

\[
f \left(x _ {1}\right) \sin \left(x _ {3} - x _ {2}\right) + f \left(x _ {2}\right) \sin \left(x _ {1} - x _ {3}\right) + f \left(x _ {3}\right) \sin \left(x _ {2} - x _ {1}\right) \geqslant 0.\tag{1}
\]

Show that:

\begin{enumerate}
\def\labelenumi{\alph{enumi})}
\tightlist
\item
  \(f\) is continuous at every point of I and has a finite right and a finite left derivative at every point of I; further,
\end{enumerate}

\[
f (x) \cos (x - y) - f _ {d} ^ {\prime} (x) \sin (x - y) \leqslant f (y)\tag{2}
\]

for every pair of points \(x, y\) of I such that \(|x - y| \leqslant \pi\) ; there is also an analogue of (2) where one replaces \(f_d'\) by \(f_g'\) ; finally, one has \(f_g'(x) \leqslant f_d'(x)\) for every \(x \in I\) . (To show (2) let \(x_2\) approach \(x_1\) in (1), keeping \(x_3\) fixed, and so obtain an upper bound for \(\limsup_{x_2 \to x_1} \frac{f(x_2) - f(x_1)}{x_2 - x_1}\) ; then let \(x_3\) approach \(x_1\) in the inequality so obtained; deduce from this the existence of \(f_d'(x)\) and the inequality (2) for \(y > x\) ; proceed similarly for the other parts.)

\begin{enumerate}
\def\labelenumi{\alph{enumi})}
\setcounter{enumi}{1}
\item
  Conversely, if (2) holds for every pair of points \(x, y\) of I such that \(|x - y| \leqslant \pi\) , then \(f\) satisfies (1) on I (consider the difference \(\frac{f(x_2)}{\sin(x_3 - x_2)} - \frac{f(x_1)}{\sin(x_3 - x_1)}\) ).
\item
  If \(f\) admits a second derivative on \(I\) , then (1) is equivalent to the condition
\end{enumerate}

\[
f (x) + f ^ {\prime \prime} (x) \geqslant 0\tag{3}
\]

for every \(x\in \mathrm{I}\)

*Interpret these results, considering the plane curve defined by \(x = \frac{1}{f(t)} \cos t\) , \(y = \frac{1}{f(t)} \sin t\) (``convexity with respect to the origin'').*

\begin{enumerate}
\def\labelenumi{\arabic{enumi})}
\setcounter{enumi}{17}
\tightlist
\item
  \begin{enumerate}
  \def\labelenumii{\alph{enumii})}
  \tightlist
  \item
    Show that in the vector space of maps from \(\mathbf{R}\) into \(\mathbf{C}\) , the distinct functions of the form \(x^n e^{\alpha x}\) ( \(n\) an integer, \(\alpha\) arbitrary complex) form a free system (argue by contradiction, considering a relation among these functions having the least possible number of coefficients \(\neq 0\) , and differentiating this relation).
  \end{enumerate}
\end{enumerate}

\begin{enumerate}
\def\labelenumi{\alph{enumi})}
\setcounter{enumi}{1}
\tightlist
\item
  Let \(f(\mathrm{X}_1, \mathrm{X}_2, \ldots, \mathrm{X}_m)\) be a polynomial in \(m\) indeterminates, with complex coefficients, such that when one substitutes the function \(\cos\left(\sum_{k=1}^{n} p_{jk} x_k\right)\) for \(\mathrm{X}_j\) for \(1 \leqslant j \leqslant r\) , and the function \(\sin\left(\sum_{k=1}^{n} p_{jk} x_k\right)\) for \(r + 1 \leqslant j \leqslant m\) , where the \(p_{jk}\) are real and the \(x_k\) are real, one obtains an identically zero function of the \(x_k\) ; show that the same identity holds when one gives the \(x_k\) arbitrary complex values (use \(a\) ).
\end{enumerate}

¶ 19) One knows (A, IX, §10, exerc. 2) that in the complex plane \(\mathbf{C}^2\) the group A of angles of directed lines is isomorphic to the orthogonal group \(\mathbf{O}_2(\mathbf{C})\) , the canonical isomorphism associating the angle \(\theta\) with rotation through the angle \(\theta\) ; one transports the topology of \(\mathbf{O}_2(\mathbf{C})\) (considered as a subspace of the space \(\mathbf{M}_2(\mathbf{C})\) of matrices of order 2 over \(\mathbf{C}\) ) to the group A by this isomorphism, so making A a locally compact topological group;

further, the map \(\theta \mapsto \cos \theta + i\sin \theta\) is an isomorphism of the topological group A onto the multiplicative group \(\mathbf{C}^*\) of complex numbers \(\neq 0\) , the inverse isomorphism being defined by the formulae \(\cos \theta = \frac{1}{2}(z + 1 / z)\) , \(\sin \theta = \frac{1}{2i} (z - 1 / z)\) . Deduce from these relations that every continuous homomorphism \(z \mapsto \varphi(z)\) of the additive group C onto A such that the complex functions \(\cos \varphi(z)\) , \(\sin \varphi(z)\) ) are differentiable on C, is defined by the relations \(\cos \varphi(z) = \cos az\) , \(\sin \varphi(z) = \sin az\) ( \(a\) a complex number).

\begin{enumerate}
\def\labelenumi{\arabic{enumi})}
\setcounter{enumi}{19}
\item
  Let D be the subset of C which is the union of the set defined by \(-\pi < \mathcal{R}(z) \leqslant \pi\) , \(\mathcal{I}(z) > 0\) , and the segment \(\mathcal{I}(z) = 0\) , \(0 \leqslant \mathcal{R}(z) \leqslant \pi\) . Show that the restriction of the function \(\cos z\) to D is a bijection of D onto C; the restriction of \(\cos z\) to the interior of D is a homeomorphism of this open set onto the complement, in C, of the half-line \(y = 0\) , \(x \leqslant 1\) .
\item
  Let \(f\) and \(g\) be two relatively prime polynomials (with complex coefficients), the degree of \(f\) being strictly less than that of \(g\) . Let \(p\) be the gcd of \(g\) and its derivative \(g'\) , and let \(q\) be the quotient of \(g\) by \(p\) ; show that there are two uniquely determined polynomials \(u, v\) whose degrees are, respectively, strictly less than those of \(p\) and \(q\) , such that
\end{enumerate}

\[
{\frac {f}{g}} = \mathrm{D} \left({\frac {u}{p}}\right) + {\frac {v}{q}}.
\]

Deduce that the coefficients of \(u\) and \(v\) belong to the smallest field (over \(\mathbf{Q}\) ) containing the coefficients of \(f\) and \(g\) and contained in \(\mathbf{C}\) .

\begin{enumerate}
\def\labelenumi{\arabic{enumi})}
\setcounter{enumi}{21}
\tightlist
\item
  Let \(f\) and \(g\) be two relatively prime polynomials (with complex coefficients), the degree of \(f\) being strictly less than that of \(g\) . Let \(K\) be a subfield of \(C\) containing the coefficients of \(f\) and \(g\) , and such that \(g\) is irreducible over \(K\) . For there to be a primitive of \(f / g\) of the form \(\sum_{i} a_i \log u_i\) , where the \(a_i\) are constants belonging to \(K\) and the \(u_i\)
\end{enumerate}

are irreducible polynomials over K, it is necessary and sufficient that \(f = cg'\) , where c is a constant in K.

\begin{enumerate}
\def\labelenumi{\arabic{enumi})}
\setcounter{enumi}{22}
\item
  If \(f(x, y)\) is an arbitrary polynomial in x, y, with complex coefficients, show that the evaluation of a primitive of \(f(x, \log x)\) and of \(f(x, \operatorname{Arc} \sin x)\) reduces to the evaluation of a primitive of a rational function.
\item
  Show that one can reduce the evaluation of a primitive of \((ax + b)^{p} x^{q}\) (p and q rational) to the evaluation of a primitive of a rational function when one of the numbers p, q, \(p + q\) is an integer (positive or negative).
\end{enumerate}

* 25) The meromorphic functions on an open disc \(\Delta\) in \(\mathbf{C}\) form a field \(\mathrm{M}(\Delta)\) . A subfield F of \(\mathrm{M}(\Delta)\) such that \(u \in \mathrm{F}\) implies \(\mathrm{D}u \in \mathrm{F}\) is called a differential subfield of \(\mathrm{M}(\Delta)\) .

\begin{enumerate}
\def\labelenumi{\alph{enumi})}
\tightlist
\item
  For every polynomial \(\mathrm{P}(\mathrm{X})=\mathrm{X}^{m}+a_{1}\mathrm{X}^{m-1}+\cdots+a_{m}\) with coefficients in \(\mathrm{M}(\Delta)\) show that there is an open disc \(\Delta_{1}\subset\Delta\) (not having the same centre as \(\Delta\) in general) and a function \(f\in\mathrm{M}(\Delta_{1})\) satisfying
\end{enumerate}

\[
\Big (f (z) \Big) ^ {m} + a _ {1} (z) \Big (f (z) \Big) ^ {m - 1} + \dots + a _ {m} (z) = 0
\]

at all the points \(z \in \Delta_1\) where \(f\) and the \(a_j\) are holomorphic. If \(F\) is a subfield of \(\mathbf{M}(\Delta)\) containing the \(a_j\) one says that the subfield of \(\mathbf{M}(\Delta_1)\) generated by \(f\) and the restrictions to \(\Delta_{1}\) of the functions of F is the field \(\mathrm{F}(f)\) obtained by adjoining the root f of P to F. This abuse of language does not cause confusion because the restriction map \(g \mapsto g|\Delta_{1}\) of F onto the subfield \(\mathrm{M}(\Delta_{1})\) is injective. If F is differential, so is \(\mathrm{F}(f)\) .

\begin{enumerate}
\def\labelenumi{\alph{enumi})}
\setcounter{enumi}{1}
\item
  For every function \(a \in \mathbf{M}(\Delta)\) show that there is an open disc \(\Delta_{2} \subset \Delta\) such that there is a function \(g \in \mathbf{M}(\Delta_{2})\) satisfying the relation \(g'(z) = a(z)\) at every point \(z \in \Delta_{2}\) where g and a are holomorphic. If F is a subfield of \(\mathbf{M}(\Delta)\) containing a one says that the subfield \(\mathrm{F}(g)\) of \(\mathbf{M}(\Delta_{2})\) generated by g and the restrictions to \(\Delta_{2}\) of the functions of F is the field obtained by adjoining the primitive \(\int a dz\) to F. If F is differential, so is \(\mathrm{F}(g)\) .
\item
  For every function \(b \in \mathrm{M}(\Delta)\) show that there is an open disc \(\Delta_{3} \subset \Delta\) such that there is an \(h \in \mathrm{M}(\Delta_{3})\) satisfying the relation \(h'(z) = b(z)h(z)\) at every point \(z \in \Delta_{3}\) where h and b are holomorphic and \(\neq 0\) . If F is a subfield of \(\mathrm{M}(\Delta)\) containing b, one says that the subfield \(\mathrm{F}(h)\) of \(\mathrm{M}(\Delta_{3})\) generated by h and the restrictions to \(\Delta_{3}\) of the functions of F is the field obtained by adjoining the exponential of the primitive \(\exp\left(\int bdz\right)\) to F. If F is differential, so is \(\mathrm{F}(h)\) .
\item
  For every subfield F of M( \(\Delta\) ) and every polynomial P(X) = X \(^{m}\) + a \(_{1}\) X \(^{m-1}\) + \(\cdots\) + a \(_{m}\) with coefficients in F there is a field K obtained by adjoining to F successively the roots of polynomials such that K is a Galois extension of F and that one has P(X) = (X - c \(_{1}\) ) \ldots{} (X - c \(_{m}\) ), where the c \(_{i}\) are meromorphic functions (on a suitable disc) belonging to K. If F is differential one has ( \(\sigma.g\) )' = \(\sigma.g'\) for every g ∈ K and every element \(\sigma\) of the Galois group of K over F.*
\end{enumerate}

* 26) a) Let \(\mathrm{F} \subset \mathrm{M}(\Delta)\) be a differential field and let \(\mathbf{K}\) be a finite Galois extension of \(\mathbf{F}\) , a subfield of a \(\mathrm{M}(\Delta_1)\) . Let \(t\) be a primitive or an exponential of a primitive of a function of \(\mathbf{F}\) ; show that if \(t\) is transcendental over \(\mathbf{F}\) there is no function \(u \in \mathbf{K}\) such that \(t' = u'\) . (Consider separately the two cases where \(t' = a \in \mathbf{F}\) or \(t' = bt\) with \(b \in \mathbf{F}\) ; obtain a contradiction by considering the transforms \(\sigma.u\) of \(u\) by the Galois group of \(\mathbf{K}\) over \(\mathbf{F}\) and their derivatives \(\sigma.u'\) ; in the first case show that one would have \(t' = c'\) for some \(c \in \mathbf{F}\) , and in the second case, putting \(\mathrm{N} = [\mathrm{K}: \mathrm{F}]\) , one would have \((t^{N}/c)' = 0\) for some \(c \in \mathrm{F}\) .)

\begin{enumerate}
\def\labelenumi{\alph{enumi})}
\setcounter{enumi}{1}
\tightlist
\item
  Suppose that t is transcendental over F and that \(t' = bt\) with \(b \in F\) ; show that there does not exist any \(c \neq 0\) in K such that \((ct^{m})' = 0\) (same method).*
\end{enumerate}

* 27) Let \(\mathrm{F} \subset \mathrm{M}(\Delta)\) be a differential field containing \(\mathbf{C}\) , and \(t\) a primitive or an exponential of a primitive of a function in \(\mathrm{F}\) , and suppose that \(t\) is transcendental over \(\mathrm{F}\) . Further, let \(c_1, \ldots, c_n\) be elements of \(\mathrm{F}\) which are linearly independent over the field of rational numbers \(\mathbf{Q}\) ; so if \(u_1, \ldots, u_n\) and \(v\) are elements of the field \(\mathrm{F}(t)\) and if

\[
\sum_ {j = 1} ^ {n} c _ {j} \frac {u _ {j} ^ {\prime}}{u _ {j}} + v ^ {\prime}
\]

belongs to the ring F{[}t{]}, then necessarily \(v \in F[t]\) . Further, if \(t \in F\) , one necessarily has \(u_{j} \in F\) for \(1 \leqslant j \leqslant n\) ; if \(t' / t \in F\) there exists for each j an integer \(v_{j} \geqslant 0\) such that \(u_{j} / t^{v_{j}} \in F\) . (Decompose the \(u_{j}' / u_{j}\) and v into simple elements in a suitable Galois extension of F and use exerc. 26).*

* 28) a) Let \(\mathrm{F} \subset \mathrm{M}(\Delta)\) be a differential field. One says that a field \(\mathrm{F}' \supset \mathrm{F}\) is an elementary extension of \(\mathrm{F}\) if there is a finite sequence

\[
\mathrm{F} = \mathrm{F} _ {0} \subset \mathrm{F} _ {1} \subset \mathrm{F} _ {2} \subset \dots \subset \mathrm{F} _ {n - 1} \subset \mathrm{F} _ {n} = \mathrm{F} ^ {\prime}\tag{1}
\]

such that for every \(j \leqslant n - 1\) one has \(\mathrm{F}_{j + 1} = \mathrm{F}_j(t_j)\) , where one has either \(t_j' = a_j' / a_j\) for an \(a_j \neq 0\) in \(\mathrm{F}_j\) (so that one can write \(t_j = \log a_j\) ), or \(t_j' / t_j = a_j'\) for an \(a_j \in \mathrm{F}_j\) (so that \(t_j = \exp(a_j)\) ). An elementary function is a function that belongs to an elementary extension of \(\mathbf{C}(z)\) (the field of rational functions over \(\mathbf{C}\) ).

For example, the functions

\[
\operatorname{Arc} \sin z, \quad (\log z) ^ {\log z}, \quad \left(1 - \exp \left(\frac {1}{\exp \left(\frac {1}{z}\right) - 1}\right)\right) ^ {\alpha} \quad (\alpha \in \mathbf {C})
\]

are elementary functions.

\begin{enumerate}
\def\labelenumi{\alph{enumi})}
\setcounter{enumi}{1}
\tightlist
\item
  Let \(a \in \mathbf{F}\) be such that there are elements \(u_{j}\) ( \(1 \leqslant j \leqslant m\) ) and \(v\) in an elementary extension \(\mathrm{F}'\) of \(\mathrm{F}\) , and constants \(\gamma_{j} \in \mathbf{C}\) ( \(1 \leqslant j \leqslant m\) ) such that
\end{enumerate}

\[
a = \sum_ {j = 1} ^ {m} \gamma_ {j} \frac {u _ {j} ^ {\prime}}{u _ {j}} + v ^ {\prime}.\tag{*}
\]

Show that then there are elements \(f_{j}\) ( \(1 \leqslant j \leqslant p\) ) and \(g\) in \(F\) and constants \(\beta_{j} \in \mathbf{C}\) ( \(1 \leqslant j \leqslant p\) ) such that

\[
a = \sum_ {j = 1} ^ {p} \beta_ {j} \frac {f _ {j} ^ {\prime}}{f _ {j}} + g ^ {\prime}.\tag{**}
\]

(Reduce, by induction on the length \(n\) of the sequence (1), to the case where \(\mathrm{F}' = \mathrm{F}(t)\) . By modifying the \(u_{j}\) first show that one can assume the \(\gamma_{j}\) to be linearly independent over \(\mathbf{Q}\) . When \(t\) is transcendental over \(\mathrm{F}\) use exerc. 27: if \(t' = s'/s\) with \(s \in \mathrm{F}\) one must have \(u_{j} \in \mathrm{F}\) and \(c \in \mathrm{F}[t]\) ; by using (*) and arguing by contradiction show that one must have \(v = \alpha t + b\) with \(\alpha \in \mathbf{C}\) and \(b \in \mathrm{F}\) . If \(t'/t = r'\) with \(r \in \mathrm{F}\) show that by replacing \(v\) by \(v + k.r\) for a suitable integer \(k \in \mathbf{Z}\) , one can assume that \(u_{j} \in \mathrm{F}\) , \(v \in \mathrm{F}[t]\) ; show by contradiction that \(v\) is of degree 0, so \(v \in \mathrm{F}\) . Finally, if \(t\) is algebraic over \(\mathrm{F}\) , embed \(\mathrm{F}'\) in a Galois extension \(\mathrm{K}\) of \(\mathrm{F}\) and consider the transforms of (*) by the Galois group of \(\mathrm{K}\) over \(\mathrm{F}.)\)*

* 29) Let \(f\) and \(g\) be two rational functions of \(z\) (elements of \(\mathbf{C}(z)\) ). Show that, for the primitive \(\int f(z)e^{g(z)}dz\) to be an elementary function (exerc. 28) it is necessary and sufficient that there exists a rational function \(r\in \mathbf{C}(z)\) such that \(f = r' + rg'\) . (Put \(t = e^g\) and consider the differential field \(F = C(z,t)\) , which is an elementary extension of \(\mathbf{C}(z)\) since \(t' / t = g'\) . Show first that \(t\) is transcendental over \(\mathbf{C}(z)\) ; arguing by contradiction, consider a Galois extension \(K\) of \(\mathbf{C}(z)\) containing \(t\) , and the transforms of the equation \(t' / t = g'\) by the Galois group of \(K\) ; one will obtain an equation \(g' = u' / u\) with \(u\in \mathbf{C}(z)\) , and it is impossible for \(g'\) to have other than simple poles in \(\mathbf{C}\) . Applying exerc. 28 one has \(ft = \sum_{j}\gamma_{j}\frac{u_{j}'}{u_{j}} +v'\) , with \(\gamma_{j}\in \mathbf{C}\) , and the \(u_{j}\) and \(v\) in \(F\) ; one can reduce to the case

where the \(\gamma_{j}\) are linearly independent over Q. Then, applying exerc. 27 to the extension \(\mathbf{C}(z)(t)\) , show that one has \(ft = v' + h\) , with \(h \in \mathbf{C}(z)\) and \(v \in \mathbf{C}(z)[t]\) . Conclude that v is necessarily of degree 1 in t, and consequently f is equal to the coefficient of t in \(v'\) .

Deduce from this that the primitives \(\int e^{z^2}dz\) and \(\int e^z dz / z\) are not elementary functions (a theorem of Liouville), by examining in the relation \(f = r' - rg'\) the behaviour of the two sides at a neighbourhood of a pole of \(r\).*

\begin{enumerate}
\def\labelenumi{\arabic{enumi})}
\setcounter{enumi}{29}
\tightlist
\item
  \begin{enumerate}
  \def\labelenumii{\alph{enumii})}
  \tightlist
  \item
    If \(m\) and \(n\) are two integers such that \(0 < m < n\) show that
  \end{enumerate}
\end{enumerate}

\[
\int_ {0} ^ {+ \infty} \frac {x ^ {m - 1}}{1 + x ^ {n}} d x = \frac {\pi}{n \sin \frac {m \pi}{n}}.
\]

\begin{enumerate}
\def\labelenumi{\alph{enumi})}
\setcounter{enumi}{1}
\tightlist
\item
  Show that for \(0 < a < 1\) the integral \(\int_0^{+\infty} \frac{x^{a-1}}{1 + x} dx\) is uniformly convergent for \(a\) varying in a compact interval, and deduce from \(a\) ) that
\end{enumerate}

\[
\int_ {0} ^ {+ \infty} \frac {x ^ {a - 1}}{1 + x} d x = \frac {\pi}{\sin a \pi}.
\]

\begin{enumerate}
\def\labelenumi{\arabic{enumi})}
\setcounter{enumi}{30}
\tightlist
\item
  If \(I_{m,n}\) is a primitive of \(\sin^{m}x\cos^{n}x\) (m and n are arbitrary real numbers), show that
\end{enumerate}

\[
\mathrm{I} _ {m + 2, n} = - \frac {\sin^ {m + 1} x \cos^ {n + 1} x}{m + n + 2} + \frac {m + 1}{m + n + 2} \mathrm{I} _ {m, n}
\]

is a primitive of \(\sin^{m+2}x\cos^{n}x\) if \(m+n+2\neq0\) .

With the aid of this formula recover the formula (29) of III, p.~100 and prove the formula

\[
\int_ {0} ^ {\pi / 2} \sin^ {2 n + 1} x d x = \frac {2 . 4 . 6 \dots 2 n}{1 . 3 . 5 \dots (2 n + 1)} \quad (n \text {   an   integer   } \geqslant 0).
\]

\begin{enumerate}
\def\labelenumi{\arabic{enumi})}
\setcounter{enumi}{31}
\tightlist
\item
  Prove Wallis' Formula
\end{enumerate}

\[
\lim _ {n \rightarrow \infty} \frac {1}{\sqrt {n}}. \frac {2 . 4 . 6 \dots 2 n}{1 . 3 . 5 \dots (2 n - 1)} = \sqrt {\pi}
\]

by using exerc. 31 and the inequality \(\sin^{n+1}x \leqslant \sin^{n}x\) for \(0 \leqslant x \leqslant \pi/2\) .

\begin{enumerate}
\def\labelenumi{\arabic{enumi})}
\setcounter{enumi}{32}
\tightlist
\item
  \begin{enumerate}
  \def\labelenumii{\alph{enumii})}
  \tightlist
  \item
    Calculate the integrals
  \end{enumerate}
\end{enumerate}

\[
\int_ {0} ^ {1} \left(1 - x ^ {2}\right) ^ {n} d x, \quad \int_ {0} ^ {+ \infty} \frac {d x}{\left(1 + x ^ {2}\right) ^ {n}}
\]

for n an integer \textgreater{} 0, with the help of exerc. 31.

\begin{enumerate}
\def\labelenumi{\alph{enumi})}
\setcounter{enumi}{1}
\tightlist
\item
  Show that one has
\end{enumerate}

\[
\begin{array}{l l} 1 - x ^ {2} \leqslant e ^ {- x ^ {2}} & \text { for } 0 \leqslant x \leqslant 1 \\ e ^ {- x ^ {2}} \leqslant \frac {1}{1 + x ^ {2}} & \text { for } x \geqslant 0. \end{array}
\]

\begin{enumerate}
\def\labelenumi{\alph{enumi})}
\setcounter{enumi}{2}
\tightlist
\item
  Deduce from \(a\) and \(b\) ) and from Wallis' formula (exerc. 32) that
\end{enumerate}

\[
\int_ {0} ^ {+ \infty} e ^ {- x ^ {2}} d x = \frac {\sqrt {\pi}}{2}.
\]

\begin{enumerate}
\def\labelenumi{\arabic{enumi})}
\setcounter{enumi}{33}
\tightlist
\item
  \begin{enumerate}
  \def\labelenumii{\alph{enumii})}
  \tightlist
  \item
    Show that for \(\alpha > 0\) the derivative of
  \end{enumerate}
\end{enumerate}

\[
\mathrm{I} (\alpha) = \int_ {0} ^ {+ \infty} e ^ {- \alpha x} \frac {\sin x}{x} d x
\]

is equal to \(-\int_{0}^{+\infty}e^{-\alpha x}\sin x dx\) .

\begin{enumerate}
\def\labelenumi{\alph{enumi})}
\setcounter{enumi}{1}
\tightlist
\item
  Deduce that \(\int_0^{+\infty}\frac{\sin x}{x} dx = \frac{\pi}{2}.\)
\end{enumerate}

\begin{enumerate}
\def\labelenumi{\arabic{enumi})}
\setcounter{enumi}{34}
\tightlist
\item
  By differentiating with respect to the parameter and using exerc. 33 c), prove the formulae
\end{enumerate}

\[
\begin{array}{l} \int_ {0} ^ {+ \infty} e ^ {- x ^ {2}} \cos \alpha x d x = \frac {\sqrt {\pi}}{2} e ^ {- \alpha^ {2}}, \quad \int_ {0} ^ {+ \infty} \frac {1 - e ^ {- \alpha x ^ {2}}}{x ^ {2}} d x = \sqrt {\pi \alpha} \\ \int_ {0} ^ {+ \infty} \exp \left(- x ^ {2} - \frac {\alpha^ {2}}{x ^ {2}}\right) d x = \frac {\sqrt {\pi}}{2} e ^ {- 2 \alpha} \quad (\alpha > 0). \end{array}
\]

\begin{enumerate}
\def\labelenumi{\arabic{enumi})}
\setcounter{enumi}{35}
\tightlist
\item
  Deduce from exerc. 33 c) above, and from II, p.~88, exerc. 9, that
\end{enumerate}

\[
\int_ {0} ^ {+ \infty} \sin x ^ {2} d x = \frac {1}{2} \sqrt {\frac {\pi}{2}}.
\]

\begin{enumerate}
\def\labelenumi{\arabic{enumi})}
\setcounter{enumi}{36}
\tightlist
\item
  Let \(\mathbf{f}\) be a regulated vector function on \(]0, 1[\) , such that the integral \(\int_0^\pi \mathbf{f}(\sin x) dx\) is convergent. Show that the integral \(\int_0^\pi x \mathbf{f}(\sin x) dx\) is convergent and that
\end{enumerate}

\[
\int_ {0} ^ {\pi} x \mathbf {f} (\sin x) d x = \frac {\pi}{2} \int_ {0} ^ {\pi} \mathbf {f} (\sin x) d x.
\]

\begin{enumerate}
\def\labelenumi{\arabic{enumi})}
\setcounter{enumi}{37}
\item
  Let \(\mathbf{f}\) be a regulated vector function for \(x \geqslant 0\) , continuous at the point \(x = 0\) , and such that the integral \(\int_{a}^{+\infty} \mathbf{f}(x) dx / x\) converges for \(a > 0\) . Show that for \(a > 0\) and \(b > 0\) the integral \(\int_{0}^{+\infty} \frac{\mathbf{f}(ax) - \mathbf{f}(bx)}{x} dx\) converges and is equal to \(\mathbf{f}(0) \log a / b\) .
\item
  Let \(m\) be a convex function on \([0, +\infty[\) such that \(m(0) = 0\) , and \(p\) a number such that \(-1 < p < +\infty\) . Show that the integral \(\int_0^{+\infty} x^p \exp(-m_r'(x)) dx\) is convergent, as is the integral \(\int_0^{+\infty} x^p \exp(-m(x)/x) dx\) , and that
\end{enumerate}

\[
\int_ {0} ^ {+ \infty} x ^ {p} \exp (- m (x) / x) d x \leqslant e ^ {p + 1} \int_ {0} ^ {+ \infty} x ^ {p} \exp (- m _ {r} ^ {\prime} (x)) d x.
\]

(For \(k > 1\) and \(A > 0\) note that \(m(kx) \geqslant m(x) + (k - 1)x m_r'(x)\) , and deduce the inequality

\[
k ^ {- p - 1} \int_ {0} ^ {k A} x ^ {p} \exp (- m (x) / x) d x \leqslant \int_ {0} ^ {A} x ^ {p} \exp \left(- \frac {m (x)}{k x} - \frac {k - 1}{k} m _ {r} ^ {\prime} (x)\right) d x.
\]

Estimate the second integral with the help of Hölder's Inequality (III, p.~115, exerc. 3) then let A tend to \(+\infty\) and \(k\) tend to 1.)

\setcounter{exercisegroup}{1}
\refstepcounter{exercisegroup}
\section*{Exercises on \S\arabic{exercisegroup}}
\phantomsection
\addcontentsline{toc}{section}{Exercises on \protect\S\arabic{exercisegroup}}

\begin{enumerate}
\def\labelenumi{\arabic{enumi})}
\tightlist
\item
  Let \(\mathbf{f}\) be a vector function that is \(n\) times differentiable on an interval \(\mathrm{I} \subset \mathbf{R}\) . Prove the formula
\end{enumerate}

\[
\mathrm{D} ^ {n} \mathbf {f} \left(e ^ {x}\right) = \sum_ {m = 1} ^ {n} \frac {a _ {m}}{m !} e ^ {m x} \mathbf {f} ^ {(m)} \left(e ^ {x}\right)
\]

at every point x such that \(e^{x} \in I\) , where the coefficient \(a_{m}\) can be expressed as

\[
a _ {m} = \sum_ {p = 0} ^ {m} (- 1) ^ {p} \binom {m} {p} (m - p) ^ {n}
\]

(method of I, p.~41, exerc. 7, using the Taylor expansion for \(e^{x}\) ).

\begin{enumerate}
\def\labelenumi{\arabic{enumi})}
\setcounter{enumi}{1}
\item
  Let \(f\) be a real function that is \(n\) times differentiable at a point \(x\) , and \(\mathbf{g}\) a vector function \(n\) times differentiable at the point \(f(x)\) . If one puts \(\mathrm{D}^n(\mathbf{g}(f(x))) = \sum_{k=1}^{n} \mathbf{g}^{(k)}(f(x)) u_k(x)\) , then \(u_k\) depends on the function \(f\) alone; deduce from this that \(u_k(x)\) is the coefficient of \(t^k\) in the expansion of \(e^{-tf(x)}\mathrm{D}^n(e^{tf(x)})\) (in terms of \(t\) ).
\item
  For every real \(x > 0\) and every complex \(m = \mu + iv\) one puts \(x^{m} = e^{m\log x}\) ; show that the formula (19) of III, p.~108 remains valid for \(m\) complex and \(x > -1\) and that the remainder \(r_{n}(x)\) in this formula satisfies the inequalities
\end{enumerate}

\[
\begin{array}{l l} | r _ {n} (x) | \leqslant \left| \binom {m - 1} {n} \frac {m}{\mu} x ^ {n} \Big [ (1 + x) ^ {\mu} - 1 \Big ] \right| & \text { if } \mu \neq 0 \\ | r _ {n} (x) | \leqslant \left| \binom {m - 1} {n} m x ^ {n} \log (1 + x) \right| & \text { if } \mu = 0. \end{array}
\]

Generalize the study of the convergence of the binomial series to the case where m is complex.

\begin{enumerate}
\def\labelenumi{\arabic{enumi})}
\setcounter{enumi}{3}
\tightlist
\item
  For every real \(x\) and every number \(p > 1\) prove the inequality
\end{enumerate}

\[
\begin{array}{r l} | 1 + x | ^ {p} \leqslant 1 + p x + \frac {p (p - 1)}{2} x ^ {2} + \dots + \frac {p (p - 1) \dots (p - m + 2)}{(m - 1) !} x ^ {m - 1} \\ & + \frac {p (p - 1) \dots (p - m + 1)}{m !} | x | ^ {m} + h _ {p} | x | ^ {p} \end{array}
\]

where one puts \(m = [p]\) (the integer part of \(p\) ) and

\[
h _ {p} = \frac {p (p - 1) \dots (p - m + 1)}{(m - 1) !} \int_ {0} ^ {1} z ^ {p - m} (1 - z) ^ {m - 1} d z.
\]

¶ 5) Show that \(\pi\) is irrational, in the following way: if one had \(\pi = p/q\) ( \(p\) and \(q\) integers), then on putting \(f(x) = (x(\pi - x))^n / n!\) , the integral \(q^n \int_0^\pi f(x) \sin x dx\) would be an integer \(>0\) (use the formula for integration by parts of order \(n + 1\) ); but show that on the other hand \(q^n \int_0^\pi f(x) \sin x dx\) tends to 0 as \(n\) tends to \(+\infty\) .

\begin{enumerate}
\def\labelenumi{\arabic{enumi})}
\setcounter{enumi}{5}
\tightlist
\item
  Show that on the interval \([-1,+1]\) the function \(|x|\) is the uniform of limit of polynomials, by remarking that \(|x|=(1-(1-x^{2}))^{1/2}\) and using the binomial series. Deduce another proof of the Weierstrass th. from this (cf.~II, p.~83, exerc. 20).
\end{enumerate}

¶ 7) Let \(p\) be a prime number and \(\mathbf{Q}_p\) the field of \(p\) -adic numbers (Gen.~Top., III, p.~322 and 323, exercs. 23 to 25), let \(\mathbf{Z}_p\) be the ring of \(p\) -adic integers, and \(\mathfrak{p}\) the principal ideal \((p)\) in \(\mathbf{Z}_p\) .

\begin{enumerate}
\def\labelenumi{\alph{enumi})}
\tightlist
\item
  Let \(a = 1 + pb\) , where \(b \in \mathbf{Z}_p\) is an element of the multiplicative group \(1 + \mathfrak{p}\) ; show that when the rational integer \(m\) increases indefinitely the \(p\) -adic number \(\frac{(1 + pb)^{p^m} - 1}{p^m}\) tends to a limit equal to the sum of the convergent series
\end{enumerate}

\[
\frac {p b}{1} - \frac {p ^ {2} b ^ {2}}{2} + \dots + (- 1) ^ {n - 1} \frac {p ^ {n} b ^ {n}}{n} + \dots
\]

One denotes this limit by \(\log a\) .

\begin{enumerate}
\def\labelenumi{\alph{enumi})}
\setcounter{enumi}{1}
\item
  Show that when the \(p\) -adic number \(x\) tends to 0 in \(\mathbf{Q}_p\) the number \(\frac{a^x - 1}{x}\) tends to \(\log a\) (use \(a\) ) and the definition of the topology of \(\mathbf{Q}_p\) .
\item
  Show that if \(p \neq 2\) one has \(\log a \equiv pb \pmod{\mathfrak{p}^2}\) , and if \(p = 2\) then \(\log a \equiv 0 \pmod{\mathfrak{p}^2}\) , and \(\log a \equiv -4b^4 \pmod{\mathfrak{p}^3}\) .
\item
  Show that if \(p \neq 2\) (resp. p = 2) then \(x \mapsto \log x\) is an isomorphism of the multiplicative topological group \(1 + p\) onto the additive topological group p (resp. \(p^{2}\) ); in particular, if \(e_{p}\) is the element of \(1 + p\) such that \(\log e_{p} = p\) (resp. \(\log e_{2} = 4\) ) then the isomorphism of \(Z_{p}\) onto \(1 + p\) , which is inverse to \(x \mapsto \frac{1}{p} \log x\) , (resp. \(x \mapsto \frac{1}{4} \log x\) ) is \(y \mapsto e_{p}^{y}\) (cf.~Gen.~Top., III, p.~323, exerc. 25).
\item
  Show that for every \(a \in 1 + p\) the continuous function \(x \mapsto a^{x}\) , defined on \(Z_{p}\) , admits a derivative equal to \(a^{x} \log a\) at every point; deduce from this that the function \(\log x\) admits a derivative equal to 1/x at every point of \(1 + p\) .
\end{enumerate}

¶ 8) a) With the notation of exerc. 7, show that the series with general term \(x^n / n!\) is convergent for every \(x \in \mathfrak{p}\) if \(p \neq 2\) , for every \(x \in \mathfrak{p}^2\) (but for no \(x \notin \mathfrak{p}^2\) ) if \(p = 2\) (determine the exponent of \(p\) in the decomposition of \(n!\) into prime factors). If \(f(x)\) is the sum of this series show that \(f\) is a continuous homomorphism of \(\mathfrak{p}\) (resp. \(\mathfrak{p}^2\) ) into \(1 + \mathfrak{p}^2\) . Deduce from this that for every \(z \in \mathbf{Z}_p\) one has \(f(pz) = e_p^z\) (resp. \(f(p^2z) = e_p^z\) ); in other words,

\[
e _ {p} = 1 + \frac {p}{1 !} + \frac {p ^ {2}}{2 !} + \dots + \frac {p ^ {n}}{n !} + \dots \quad \text { if } p \neq 2
\]

\[
e _ {2} = 1 + \frac {4}{1 !} + \frac {4 ^ {2}}{2 !} + \dots + \frac {4 ^ {n}}{n !} + \dots
\]

(use exerc. 7 e)).

\begin{enumerate}
\def\labelenumi{\alph{enumi})}
\setcounter{enumi}{1}
\tightlist
\item
  For every \(a \in 1 + \mathfrak{p}\) and every \(x \in \mathbf{Z}_p\) show that \(\log (a^x) = x\log a\) , and deduce from a) and exerc. 7 d) that
\end{enumerate}

\[
a ^ {x} = 1 + \frac {x \log a}{1 !} + \frac {x ^ {2} (\log a) ^ {2}}{2 !} + \dots + \frac {x ^ {n} (\log a) ^ {n}}{n !} + \dots
\]

\begin{enumerate}
\def\labelenumi{\alph{enumi})}
\setcounter{enumi}{2}
\tightlist
\item
  For every \(m \in Z_{p}\) show that the continuous function \(x \mapsto x^{m}\) , defined on \(1 + p\) , admits a derivative equal to \(mx^{m-1}\) (use b) and exerc. 7 e)).
\end{enumerate}

\(d\) ) Show that for every \(m\in \mathbf{Z}_p\) and every \(x\in \mathfrak{p}\) if \(p\neq 2\) ( \(x\in \mathfrak{p}^2\) if \(p = 2\) ) the series with general term \(\binom{m}{n}x^n\) is convergent and that its sum is a continuous function of \(m\) ; deduce from this that this sum is equal to \((1 + x)^m\) by remarking that \(\mathbf{Z}\) is (everywhere) dense in \(\mathbf{Z}_p\) .

¶ 9) With the notation of exerc. 7) one denotes by \(\mathbf{O}_2^+ (\mathbf{Q}_p)\) (the group of rotations of the space \(\mathbf{Q}_p^2\) ) the group of matrices of the form

\[
\left( \begin{array}{c c} x & y \\ - y & x \end{array} \right)
\]

with elements in \(\mathbf{Q}_p\) , such that \(x^2 + y^2 = 1\) , this group being endowed with the topology defined in Gen.~Top., VIII, p.~125, exerc. 2.

\begin{enumerate}
\def\labelenumi{\alph{enumi})}
\item
  Denote by \(G_{n}\) the subgroup of \(\mathbf{O}_{2}^{+}(\mathbf{Q}_{p})\) formed by the matrices such that \(t=y/(1+x)\in\mathfrak{p}^{n}\) . Show that \(G_{n}\) is a compact group, that \(G_{n}/G_{n+1}\) is isomorphic to Z/pZ, and that the only compact subgroups of \(G_{1}\) are the groups \(G_{n}\) (cf.~Gen.~Top., III, p.~323, exerc. 24).
\item
  Show that \(G_{1}\) is identical to the subgroup of matrices
\end{enumerate}

\[
\left( \begin{array}{c c} x & y \\ - y & x \end{array} \right)
\]

such that \(x^{2} + y^{2} = 1\) , \(x \in 1 + p^{2}\) and \(y \in p\) if \(p \neq 2\) ; to \(x \in 1 + p^{3}\) and \(y \in p^{2}\) if p = 2.

\begin{enumerate}
\def\labelenumi{\alph{enumi})}
\setcounter{enumi}{2}
\tightlist
\item
  Show that the series with general terms \((-1)^{n} x^{2n}/(2n)!\) and \((-1)^{n-1} x^{2n+1}/(2n+1)!\) are convergent for every \(x \in p\) if \(\neq 2\) , for every \(x \in p^{2}\) if p = 2; let \(\cos x\) and \(\sin x\) be the sums of these series. Show that the map
\end{enumerate}

\[
x \mapsto \left( \begin{array}{c c} \cos x & \sin x \\ - \sin x & \cos x \end{array} \right)
\]

is an isomorphism of the additive topological group \(\mathfrak{p}\) (resp. \(\mathfrak{p}^2\) ) onto the group \(G_1\) .

\(d\) ) If \(p\) is of the form \(4h + 1\) ( \(h\) an integer), there exists in \(\mathbf{Q}_p\) an element \(i\) such that \(i^2 = -1\) . If with every \(z \in \mathbf{Q}_p^*\) one associates the matrix

\[
\left( \begin{array}{c c} \frac {1}{2} \left(z + \frac {1}{z}\right) & \frac {1}{2 i} \left(z - \frac {1}{z}\right) \\ - \frac {1}{2 i} \left(z - \frac {1}{z}\right) & \frac {1}{2} \left(z + \frac {1}{z}\right) \end{array} \right)
\]

one defines an isomorphism of the multiplicative group \(Q_{p}^{*}\) onto the group \(\mathbf{O}_{2}^{+}(\mathbf{Q}_{p})\) ; under this isomorphism the group \(1 + p\) corresponds to \(G_{1}\) and one has \(\cos px + i \sin px = e_{p}^{ix}\) (III, p.~126, exerc. 8).

\begin{enumerate}
\def\labelenumi{\alph{enumi})}
\setcounter{enumi}{4}
\tightlist
\item
  If p is of the form \(4h + 3\) (h an integer) the elements of the matrices of \(\mathbf{O}_{2}^{+}(\mathbf{Q}_{p})\) are necessarily p-adic integers. The polynomial \(X^{2} + 1\) is then irreducible in \(Q_{p}\) ; let \(\mathbf{Q}_{p}(i)\) be the quadratic extension of \(Q_{p}\) obtained by adjoining a root i of \(X^{2} + 1\) ; one endows \(\mathbf{Q}_{p}(i)\) with the topology defined in Gen.~Top., VIII, p.~127, exerc. 2. The group \(\mathbf{O}_{2}^{+}(\mathbf{Q}_{p})\) is isomorphic to the multiplicative group N of elements of \(\mathbf{Q}_{p}(i)\) of norm 1, under the isomorphism which associates the matrix \(\begin{pmatrix} x & y \\ -y & x \end{pmatrix}\) with the element \(z = x + iy\) . Show that there are \(p + 1\) roots of the equation \(x^{p+1} = 1\) in \(\mathbf{Q}_{p}(i)\) and that they form a cyclic subgroup R of N (argue as in Gen.~Top., III, p.~323, exerc. 24; then show that there are \(p + 1\) distinct roots of the congruence \(x^{p+1} \equiv 1 \pmod{p}\) in \(\mathbf{Q}_{p}(i)\) , and, for each root a of this congruence, form the series \((a^{p2n})\) ). Deduce that the group \(\mathbf{O}_{2}^{+}(\mathbf{Q}_{p})\) is isomorphic to the product of the groups R and \(G_{1}\) .
\end{enumerate}

\(f\) ) Show that, for \(p = 2\) , the group \(\mathbf{O}_2^+ (\mathbf{Q}_p)\) is isomorphic to the product of the group \(\mathrm{G}_1\) and the cyclic group of order 4.
